% !PRINTED 543
\weberpage{543}
%<<CONT
Klassenzahl selbst, sondern den dritten Teil oder die Hälfte davon versteht.

Nun ist $\delta'$ immer gerade und von entgegengesetztem Vorzeichen wie $\delta$, und $4m=-\delta\delta'$; danach ergeben sich aus (7) die Formeln:
\weq{10}{\begin{aligned}
2\overset{k}{\Sigma}\,\chi(\delta,k)\log\frac{f(\omega)}{\sqrt[4]{2}} &= 0, & \delta&\equiv 1\ (\mathrm{mod}\ 8), \\
&= K(\delta)K(\delta')\log\varepsilon, & \delta&\equiv 5\ (\mathrm{mod}\ 8),
\end{aligned}}
worin
% [sic] The source appends incongruous residue conditions to delta>0 and delta<0 in (11); the unit is selected by the sign of delta, not by these printed residue clauses.
\weq{11}{\begin{aligned}
\varepsilon &= \frac{T+U\sqrt{\delta}}{2}, & \delta&>0\equiv 1\ (\mathrm{mod}\ 8), \\
&= \frac{T+U\sqrt{\delta'}}{2}, & \delta&<0\equiv 5\ (\mathrm{mod}\ 8)
\end{aligned}}
zu setzen ist. Die erste der Formeln (10) gilt auch noch für $\delta=1$.

Da nun für alle Klassen eines Geschlechts und für jedes $\delta$ der Charakter $\chi(\delta,k)$ denselben Wert hat, so sind durch (10) und (11) die Produkte (1) bestimmt. Denn es ist nach \S\,113
\[
\overset{\delta}{\Sigma}\,\chi(\delta,k) = 0,
\]
% [sic] The source says Hauptklasse; the genus-character sum is g for every class in the principal genus, not only the principal class.
außer wenn $k$ die Hauptklasse ist, und für diese ist die Summe gleich der Anzahl $g$ der Geschlechter. Man erhält z.\,B. für die Wurzeln $\omega$ des \gesperrt{Hauptgeschlechts}
\[
\left(\Pi\,\frac{f(\omega)}{\sqrt[4]{2}}\right)^{2g} = \Pi\,\varepsilon^{K(\delta)K(\delta')},
\]
worin $g$ die Anzahl der Geschlechter bedeutet, und das Produkt links über alle Wurzeln $\omega$ des Hauptgeschlechts, das Produkt rechts über alle Stammteiler $\delta$ von $\Delta$, die $\equiv 5\ (\mathrm{mod}\ 8)$ sind, auszudehnen ist.

% Editorial cross-reference: (10) is the evaluated form; weighting (8) and applying (7),(9) also gives the genus-product identity. The printed reference (8) is retained.
Um das Produkt der Klasseninvarianten für ein anderes als das Hauptgeschlecht zu bilden, hat man die Formel (8) vor der Summation mit $\chi(\delta,k^{-1})$ zu multiplizieren, wenn $k$ eine Klasse des betreffenden Geschlechts bedeutet.

Die Anwendung der Formel (10) verlangt die Kenntnis der Klassenzahlen positiver und negativer Diskriminanten und der Zahlen $T$, $U$. Die Klassenzahlen sind in weitem Umfange von \gesperrt{Gauss} berechnet und aus seinem Nachlaß im zweiten Bande der Werke, S.\,450 bis 476, veröffentlicht. Für die Lösungen $T$, $U$ der \gesperrt{Pell}schen Gleichung enthält \gesperrt{Legendres} \glqq Theorie des nombres\grqq{} oder auch der \glqq Canon Pellianus\grqq{} von \gesperrt{Degen} eine Tabelle.
